Trong suốt quá trình thực hiện đề tài tốt nghiệp phát triển trò chơi 2D "Dream Witch", Em đã nghiên cứu, thiết kế và cài đặt một số giải pháp kỹ thuật cốt lõi nhằm giải quyết bài toán về hiệu năng vận hành, quản lý logic trạng thái nhân vật phức tạp, cũng như kiểm soát luồng hội thoại phi tuyến. Chương này trình bày chi tiết ba giải pháp đóng góp nổi bật của đề tài.

\section{Tối ưu hóa hiệu năng bằng kỹ thuật Distance Culling tùy biến}

\subsection{Đặt vấn đề và bài toán tối ưu}
Trong thể loại game Metroidvania, môi trường thường bao gồm một lưới bản đồ liền mạch chứa số lượng lớn vật thể động (quái vật, bẫy, hoạt ảnh Spine 2D/Animator). Trong cấu trúc mặc định của Unity, tất cả các thực thể được kích hoạt trên Scene đều liên tục thực thi các hàm cập nhật logic (\texttt{Update}, \texttt{FixedUpdate}) để tính toán vật lý, phát hiện va chạm và chạy AI định tuyến (A* Pathfinding). 

Đặc biệt, Spine 2D skeletal animation của các thực thể Boss khổng lồ tiêu tốn rất nhiều CPU cho phép toán biến dạng lưới đỉnh (Mesh Deformation) và nội suy khung xương thời gian thực. Việc vận hành đồng thời toàn bộ các tài nguyên này dễ dẫn đến quá tải CPU (CPU Bottleneck), làm giảm FPS dưới mức tối thiểu (30 FPS) trên các máy tính cấu hình trung bình hoặc thấp. Vì vậy, bài toán đặt ra là cần xây dựng giải pháp kỹ thuật vô hiệu hóa chọn lọc cả logic cập nhật lẫn đồ họa của các thực thể ngoài tầm nhìn của người chơi mà vẫn bảo đảm tính nhất quán của trạng thái thế giới game.

\subsection{Giải pháp hiện thực hóa}
Để giải quyết bài toán trên, Em đề xuất giải pháp thiết kế hệ thống quản lý culling khoảng cách tùy biến dựa trên hai thành phần cốt lõi: \texttt{CullingManager} và \texttt{DistanceCullingTarget} (như đã giới thiệu sơ bộ ở Chương 5). So với Camera Frustum Culling truyền thống (chỉ dừng render nhưng vẫn chạy CPU logic cập nhật và vật lý), giải pháp này thực hiện ngắt chọn lọc cả logic cập nhật và hiển thị dựa trên khoảng cách vật lý thực tế.

Khoảng cách kiểm tra được xác định bằng công thức khoảng cách Euclid giữa vị trí người chơi $P(x_p, y_p)$ và đối tượng culling $T(x_t, y_t)$:
\begin{equation}
d = \sqrt{(x_p - x_t)^2 + (y_p - y_t)^2}
\end{equation}

Các đối tượng (\texttt{DistanceCullingTarget}) được phân loại dựa trên bán kính kích hoạt ($R_c$):
\begin{itemize}
    \item \textbf{Quái vật thường (Enemy):} $R_c = 25$ đơn vị đo lường Unity.
    \item \textbf{Boss lớn:} $R_c = 40$ đơn vị (bán kính lớn hơn để tránh ngắt AI khi Boss di chuyển ngoài khung hình).
    \item \textbf{Bẫy môi trường động:} $R_c = 20$ đơn vị.
    \item \textbf{Vật thể trang trí tĩnh có diễn họa (VFX/Deco):} $R_c = 15$ đơn vị.
\end{itemize}

Quy trình xử lý culling diễn ra như sau:
\begin{enumerate}
    \item \texttt{CullingManager} quản lý danh sách các đối tượng thông qua cấu trúc \texttt{HashSet} tối ưu để đạt độ phức tạp tìm kiếm $O(1)$.
    \item Định kỳ mỗi $0.25$ giây (chu kỳ tối ưu giúp giảm tần suất tính toán so với thực hiện từng khung hình), hệ thống duyệt qua danh sách đối tượng và tính khoảng cách $d$ tới người chơi.
    \item Nếu $d > R_c$: Đối tượng chuyển sang trạng thái \textit{Culled}, hệ thống thực hiện vô hiệu hóa:
    \begin{itemize}
        \item Component di chuyển vật lý \texttt{Rigidbody2D} và \texttt{Collider2D} (giải phóng Physics2D).
        \item Thành phần \texttt{Animator} hoặc Spine (\texttt{SkeletonAnimation}) để dừng tính toán hoạt ảnh.
        \item Luồng xử lý AI tìm đường A* (\texttt{AIPath} và \texttt{Pathfinder} components).
    \end{itemize}
    \item Nếu $d \le R_c$: Khôi phục tức thời các thành phần trên về trạng thái hoạt động ban đầu mà không gây trễ hiển thị đối với người chơi.
\end{enumerate}

\subsection{Kết quả đạt được}
Giải pháp culling tùy biến đã đem lại hiệu quả tối ưu hóa hiệu năng rõ rệt:
\begin{itemize}
    \item \textbf{FPS ổn định:} Tại các khu vực mật độ quái và bẫy dày đặc, tốc độ khung hình trung bình của trò chơi tăng từ 42 FPS lên mức 60 FPS ổn định.
    \item \textbf{Giảm tải CPU:} Thời gian CPU dành cho hoạt ảnh giảm 75\% (từ 12ms xuống dưới 3ms mỗi frame). Overhead va chạm của Physics Engine đối với các đối tượng ngoài màn hình giảm 80\%.
    \item \textbf{Trải nghiệm liền mạch:} Buffer khoảng cách nhỏ giúp kích hoạt mượt mà, loại bỏ hiện tượng các thực thể đột ngột xuất hiện (pop-in).
\end{itemize}

\section{Kiến trúc máy trạng thái phân tầng cho hệ thống nhân vật phức tạp}

\subsection{Đặt vấn đề và bài toán thiết kế}
Nhân vật chính trong game hành động đòi hỏi tập hợp hành vi di chuyển và chiến đấu đa dạng (chạy, nhảy đơn/đôi, nhảy bám tường, lướt - Dash, leo tường, combo cận chiến, né lướt, đỡ đòn). Đồng thời, cơ chế biến hình (Form Switch) thay đổi các thông số vật lý (trọng lực, tốc độ, lực nhảy) và bộ kỹ năng đi kèm.

Nếu viết tất cả logic này trong một lớp duy nhất sử dụng cấu trúc \texttt{if-else} hoặc \texttt{switch-case} lồng nhau, mã nguồn sẽ phức tạp (spaghetti code), khó bảo trì và dễ lỗi logic trạng thái điều khiển (như tấn công khi đang lướt tự do). Vì vậy, bài toán đặt ra là thiết kế một kiến trúc điều khiển tách biệt, an toàn về trạng thái và có khả năng mở rộng tốt.

\subsection{Giải pháp hiện thực hóa}
Em đã thiết kế và cài đặt kiến trúc máy trạng thái phân tầng (\textit{Hierarchical State Machine - HSM}) kết hợp mô hình hướng đối tượng. Hệ thống gồm ba lớp chính:
\begin{itemize}
    \item \texttt{PlayerController}: Lớp trung gian điều phối chứa tham chiếu vật lý, quản lý hoạt ảnh, thuộc tính nhân vật và đối tượng máy trạng thái.
    \item \texttt{PlayerStateMachine}: Quản lý trạng thái hiện tại và thực hiện chuyển dịch trạng thái qua \texttt{Initialize()} và \texttt{ChangeState()}.
    \item \texttt{PlayerState}: Lớp cơ sở trừu tượng quy định các hàm vòng đời: \texttt{Enter()} (bắt đầu), \texttt{LogicUpdate()} (cập nhật input mỗi khung hình), \texttt{PhysicsUpdate()} (cập nhật vật lý định kỳ) và \texttt{Exit()} (dọn dẹp khi thoát).
\end{itemize}

Các trạng thái cụ thể như \texttt{IdleState}, \texttt{JumpState}, \texttt{DashState}, và \texttt{WallClimbState} kế thừa từ \texttt{PlayerState}. Để giải quyết việc chuyển dạng (Form Switch) mà không làm gián đoạn máy trạng thái, Em tích hợp thêm phân hệ \texttt{PlayerFormManager} hoạt động song song.

Quy trình xử lý Form Switch diễn ra như sau:
\begin{enumerate}
    \item \texttt{PlayerFormManager} đọc dữ liệu từ ScriptableObject cấu hình tương ứng (\texttt{PlayerFormDataSO}).
    \item Thay vì khởi tạo lại toàn bộ máy trạng thái, hệ thống cập nhật trực tiếp tập hợp hoạt ảnh \texttt{Animator} và các thông số runtime (Max HP, tốc độ di chuyển, lực nhảy) trong \texttt{PlayerController}.
    \item Trạng thái hiện tại (\texttt{currentState}) nhận diện chỉ số vật lý mới và áp dụng trong hàm \texttt{PhysicsUpdate()} tiếp theo, giúp nhân vật chuyển động liên tục.
\end{enumerate}

\subsection{Kết quả đạt được}
Kiến trúc này mang lại những cải tiến quan trọng:
\begin{itemize}
    \item \textbf{Loại bỏ lỗi xung đột logic:} Logic được đóng gói độc lập trong từng lớp trạng thái. Nhân vật không thể hành động sai quy định (như cấm tấn công khi đang leo tường vì \texttt{WallClimbState} không xử lý hành vi tấn công).
    \item \textbf{Khả năng mở rộng cao:} Việc thêm trạng thái hay dạng biến hình mới chỉ yêu cầu tạo lớp kế thừa \texttt{PlayerState} hoặc file cấu hình \texttt{PlayerFormDataSO} mới mà không cần chỉnh sửa mã nguồn hiện có (tuân thủ Open/Closed Principle).
    \item \textbf{Chuyển đổi mượt mà:} Animator mới và chỉ số được cập nhật tức thời giúp nhịp độ chiến đấu liền mạch.
\end{itemize}

\section{Giải pháp tích hợp Yarn Spinner và quản lý trạng thái thế giới phi tuyến}

\subsection{Đặt vấn đề}
Trò chơi phiêu lưu yêu cầu hệ thống hội thoại rẽ nhánh và phụ thuộc chặt chẽ vào tiến trình thế giới (Quest progression) như vật phẩm đã thu thập, cổng dịch chuyển đã kích hoạt hay danh sách Boss đã hạ gục.

Việc tự xây dựng hệ thống hội thoại rẽ nhánh phức tạp từ đầu đòi hỏi nhiều thời gian lập trình giao diện, phân tích cú pháp hội thoại và dễ lỗi font tiếng Việt. Đồng thời, đồng bộ luồng hội thoại động với dữ liệu trạng thái thế giới và hệ thống lưu trữ bền vững (Save/Load) qua JSON là yêu cầu kỹ thuật quan trọng.

\subsection{Cơ chế tích hợp và đồng bộ hóa trạng thái}
Em giải quyết bài toán trên bằng cách tích hợp thư viện Yarn Spinner kết hợp giải pháp lưu trữ của \texttt{GameSaveManager}.

Giải pháp kỹ thuật bao gồm ba thành phần chính:
\begin{itemize}
    \item \textbf{Kịch bản Yarn Script:} Các cuộc hội thoại được viết trên các tệp \texttt{.yarn} độc lập. Các điều kiện rẽ nhánh sử dụng các biến số nội tại của Yarn (ví dụ: \texttt{\$has\_defeated\_golem}).
    \item \textbf{Lớp cầu nối YarnDialogueView:} Chịu trách nhiệm hiển thị UI hội thoại tiếng Việt có dấu và các lựa chọn rẽ nhánh cho người chơi.
    \item \textbf{Cơ chế đồng bộ hóa biến số (Variable Storage Synchronizer):} Khi nạp game, \texttt{GameSaveManager} đọc tệp JSON (\texttt{GameSaveData}) gồm thông tin Boss đã hạ gục (\texttt{bossDefeats}), cửa đã mở (\texttt{doors}), cổng dịch chuyển đã đi qua (\texttt{portals}) và tiến trình kỹ năng (\texttt{skillTreeProgress}) để nạp vào Yarn Spinner qua \texttt{VariableStorage.SetValue()}.
\end{itemize}

Ngược lại, khi các lựa chọn hội thoại làm thay đổi biến tiến trình trong Yarn Spinner, các sự kiện callback sẽ gọi hàm tương ứng trong \texttt{GameSaveManager} để cập nhật tệp JSON thời gian thực, bảo đảm đồng bộ tuyệt đối.

\subsection{Kết quả đạt được}
Giải pháp tích hợp mang lại các kết quả thực tiễn:
\begin{itemize}
    \item \textbf{Quy trình viết kịch bản chuyên nghiệp:} Tách biệt nội dung hội thoại ra khỏi mã nguồn C\#. Nhà thiết kế game có thể cập nhật tệp \texttt{.yarn} mà không cần biên dịch lại dự án.
    \item \textbf{Lưu giữ dữ liệu đồng bộ:} Trạng thái hội thoại và tiến trình thế giới lưu trực tiếp vào tệp JSON. Khi tải lại game, NPC hiển thị đúng lời thoại tương ứng với trạng thái nhiệm vụ trước đó.
    \item \textbf{Giao diện tiếng Việt hoàn chỉnh:} Hệ thống hoạt động ổn định, hỗ trợ hiển thị tiếng Việt mượt mà và tương thích tốt với các điều khiển đầu vào.
\end{itemize}